\section{Definitions}\label{sec:offline-definitions}

\subsection{Offline payments}

An offline payment system involves five types of participants: the central bank,
intermediaries, users, a registration authority, and an auditor.

The central bank issues tokens to users and manages the reserves of the
intermediaries. Issuance is triggered by a withdrawal request from an
intermediary and, when successful, debits the issued value from that
intermediary's reserves. The central bank also redeems tokens on a deposit
request from an intermediary, crediting the redeemed value to the requesting
intermediary's reserves.

Users hold bank deposits at intermediaries, which they can exchange for tokens.
A user withdraws a token by submitting a withdrawal request to their
intermediary, which routes the request to the central bank. A successful
withdrawal debits the withdrawn value from the user's deposits and from the
intermediary's reserves. Once a user holds a token, they can use it in offline
payments. The token's lifecycle ends with a deposit, which lets its owner claim
bank deposits in exchange for the token. A deposit succeeds only if the token
was issued by the central bank, was initiated by its owner, and had not
previously been deposited.

Intermediaries maintain users' bank deposit accounts and mediate the interaction
of users with the central bank. They also onboard users into the system, opening
accounts for them and helping them obtain long term credentials from the
registration authority. The registration authority issues these long term
credentials, which users subsequently use to authorise and authenticate
payments. The auditor is authorised to trace a fraudulent offline payment, such
as one that double spends a token, to its originator.

We instantiate each user's wallet with a mobile device and a secure element,
which together let the user perform offline payments securely and accountably.

We assume the registration authority, the central bank, and the intermediaries
are honest but curious: they follow the protocol, but may collude with other
participants to learn information about offline payments.  The registration
authority is additionally trusted to assign each user a unique identifier, to
generate a correctly formed credential that will be accepted by all other
participants, and to issue that credential only following well established KYC
processes. Users may be arbitrarily malicious: they may collude to forge tokens,
double spend them, or steal the tokens of honest users, and may collude with
other participants to undermine the privacy of honest users. Users have full
control of their mobile devices, but we assume they cannot easily tamper with
their secure elements. The auditor is fully trusted to inspect only fraudulent
payments and to deanonymise only misbehaving users. We consider static
corruption throughout: the adversary fixes the set of corrupt participants
before any protocol interaction takes place.

An offline payment scheme is defined by the following algorithms. Account
creation is an operation of the intermediary rather than of the scheme, so we do
not give it as an algorithm here, though the ideal functionality below exposes
an interface for it. Throughout, we suppress the public parameters, the long
term credential, and the counterparty's inputs from the algorithm signatures,
and name only the state that distinguishes one call from another.

\begin{description}

  \item[$(\pk_{\algo{cb}}, \ck, \pk_\auditor, \rapk) \leftarrow
\OPgen(\secparam)$ :] The setup algorithm, on input security parameter
$\secparam$, generates the public keys of the central bank, the auditor, and the
registration authority, together with a public commitment key $\ck$.

  \item[$\cred \leftarrow \OPregister(\rask)$ :] The registration algorithm, on
input the registration authority's secret key $\rask$, provisions the user's
secure element and outputs long term credential $\cred$.

  \item[$T_0 \leftarrow \OPwithdraw(v_0)$ :] The withdrawal algorithm, on input
value $v_0$, outputs token $T_0$ of that value, obtained by the user from the
central bank through their intermediary.

  \item[$p_i \leftarrow \OPexecute(v, s_i, \com_i, \algo{hist}_i)$ :] The
payment algorithm, on input a token of value $v$ with randomness $s_i$,
pseudonym $\com_i$, and history $\algo{hist}_i$, outputs payment $p_i$
transferring the token to the payee.

  \item[$\{0,1\} \leftarrow \OPdeposit(\tau_n, \algo{hist}_n)$ :] The deposit
algorithm, on input token $\tau_n$ and its history $\algo{hist}_n$, returns 1 if
the token is redeemed for commercial bank money and 0 otherwise.

  \item[$\algo{DS} \leftarrow \OPidentify(d_n, d_m)$ :] The identification
algorithm, on input two conflicting deposit requests $d_n$ and $d_m$ for the
same original token, outputs the set $\algo{DS}$ of identifiers of the users who
double spent it.

\end{description}

We define security of an offline payment scheme through an ideal functionality
that represents a trusted third party which executes the operations of the
payment system on behalf of system participants in a way that perfectly meets
the system requirements. The properties this functionality captures are stated
informally below and formalised in~\Cref{sec:ideal-func}.

\subsection{Security goals}

An offline payment system must protect against forgery, token theft, and double
spending, must conserve value, and must preserve the privacy of users against
the central bank and the intermediaries. More specifically, it must ensure:

\begin{description}

  \item[Unforgeability.] An honest payee should not be tricked, during a
payment, into accepting a token that was not issued by the central bank.
Similarly, the honest but curious central bank will not declare a deposit
successful unless it previously issued the corresponding token.

  \item[Theft prevention.] A malicious payer should not be able to convince an
honest payee to accept a token the payer does not own. Likewise, a malicious
user should not be able to successfully deposit a token they do not own.

  \item[Double spending detection and identification.] An honest user uses a
token only once, in either a payment or a deposit, after which the token is
deleted. A malicious user may instead use a token in multiple payments or
deposits. A secure offline payment system must ensure that such double spending
is swiftly detected at the time of deposit, and that the malicious user is
identified by the authorised auditor.

  \item[Conservation of value.] The total value in the system must not change
other than through issuance and redemption by the central bank. No sequence of
payments may create value, and the value credited on deposit must be the value
that was withdrawn.

  \item[Non-frameability.] An honest user should never be accused of double
spending by the authorised auditor, even if every other user, the
intermediaries, and the central bank collude.

  \item[User anonymity.] Only the intermediary of an honest user making a
withdrawal should learn their identity. During a payment, the history of the
token being transferred leaks no information about the honest users involved in
previous payments beyond what is unavoidably revealed. What is unavoidably
revealed is that the payee of the last payment is the payer of the current one,
together with a growing list of random values fixed throughout the token's
lifetime. We assume the payer and the payee know each other, and that this
knowledge is not derived from the payment transcript. The same holds for
deposits and withdrawals. The intermediary of a withdrawing or depositing user
knows their identity, but the history of a deposited token leaks nothing about
the honest users involved in payments preceding the deposit, beyond that same
list of random values and the intermediary that identifies the depositor for
accounting purposes.

  \item[Token unlinkability.] The intermediaries and the central bank should not
be able to link the withdrawal of an honest user to a deposit, beyond what is
revealed by the value of the withdrawal and the deposit, even if the
intermediaries and the central bank collude with every other user in the system.

\end{description}

\subsection{Ideal functionality $\idfu{\algo{op}}$}\label{sec:ideal-func}

We define security for the offline payment protocol through an ideal
functionality $\idfu{\algo{op}}$, which represents an incorruptible trusted
party that receives the inputs of each protocol task and returns the prescribed
outputs. We have the outputs of the honest parties and $\adv$ during an
execution of protocol $\algo{prot}$ in the real world captured in
$\textsc{Exec}_{\algo{prot},\adv}(\secparam)$, and the outputs of the honest
parties and $\simu$ interacting with $\idfu{\algo{op}}$ in the ideal world
captured by $\textsc{Exec}_{\idfu{\algo{op}},\simu}(\secparam)$. We then say
that $\algo{prot}$ securely realises $\idfu{\algo{op}}$ if for every PPT
adversary $\adv$ there exists a PPT simulator $\simu$ such that
$\textsc{Exec}_{\algo{prot},\adv}(\secparam)$ is computationally
indistinguishable from $\textsc{Exec}_{\idfu{\algo{op}},\simu}(\secparam)$. This
is a standalone simulation based definition rather than a universally composable
one, and the distinction matters: our simulator extracts from some proofs by
rewinding the adversary, which a universally composable simulator may not do. We
return to this point in~\Cref{sec:offline-security}.

$\idfu{\algo{op}}$ interacts with $\simu$ and with the system's participants.
These are the users, who act as both payers and payees, the intermediaries, the
registration authority, the central bank, and the auditor. Each user is
represented as a pair consisting of a mobile device and a secure element, so
that an individual holding multiple bank accounts is modelled as multiple such
users, one per account.  The registration authority, the central bank, and the
intermediaries are honest but curious, so $\simu$ learns their inputs and
outputs by default. We consider a static corruption model, in which $\simu$
corrupts the devices and secure elements of its choosing before they interact
with $\idfu{\algo{op}}$.  The formal description is given
in~\Cref{fig:ideal-func}. We summarise each interface below.

\paragraph{Corrupt.} $\simu$ calls this to compromise a user's mobile device or
secure element. We assume that corrupting the secure element also corrupts the
device, since compromising a secure element is significantly harder in practice.

\paragraph{Create.} An intermediary calls this to open an account for a user
with an initial value. Each user may hold only one account in the system.

\paragraph{Register.} The registration authority calls this to enrol a user who
already holds an account. A user can be registered only once.

\paragraph{Withdraw.} An intermediary calls this on behalf of one of its users,
indicating the value to withdraw. If the user's account and the intermediary's
reserves hold enough value, the user obtains a token identified by a fresh
random number, ready for use in subsequent offline payments. This identifier is
kept hidden from the intermediary and the central bank, to ensure a withdrawal
cannot later be linked to a deposit. It is chosen by $\simu$ whenever the user's
secure element is corrupt, since it is the secure element that generates it in
the protocol.  $\simu$ learns it whenever either the device or the secure
element is corrupt, since the device sees it too, and otherwise learns only the
user's identity and the withdrawn value, reflecting that the intermediary and
the central bank are honest but curious.

\paragraph{Execute.} A registered user calls this to pay another registered
user, transferring ownership of a token they hold. $\idfu{\algo{op}}$ checks
that the payee is willing to accept the payment, and that the token genuinely
belongs to the payer and has the claimed value. If so, the payee receives the
token under a fresh identifier that lets them transfer it onward in turn. A
corrupt payee may choose to accept an invalid token, but can never pass it on to
an honest user, since only genuinely valid tokens are recorded as owned.
$\simu$ learns the identities of the payer and payee only when one of their
devices or secure elements is corrupt, and otherwise learns only the value and
the token's history.

\paragraph{Deposit.} An intermediary calls this to let a registered user redeem
a token for commercial bank money. $\idfu{\algo{op}}$ checks that the token
genuinely belongs to the depositing user, has the claimed value, and has not
already been deposited: a second deposit of the same withdrawal signals double
spending and the call fails. Otherwise, the user's account and the
intermediary's reserves are credited accordingly, and by exactly the value that
was withdrawn, which is what gives conservation of value. $\simu$ learns the
deposit information by default, since the central bank and the intermediary are
honest but curious.

\paragraph{Identify.} The auditor calls this with the identifier of a withdrawn
token, to find out whether it has been double spent. If it has,
$\idfu{\algo{op}}$ identifies each double spender as the owner of the token at a
point where two deposit chains for it diverge, and returns the resulting set to
the auditor and the central bank. Settlement of the double spent value against
those users is an operation of the central bank rather than of the scheme, so
$\idfu{\algo{op}}$ reports the responsible parties and does not model the
accounting that follows.

\paragraph{How the goals are captured.} Each informal goal above corresponds to
a specific behaviour of $\idfu{\algo{op}}$. Unforgeability is the check in
$\OPdeposit$ and $\OPexecute$ that $\algo{tokens}[\vec{\rho}]$ is defined, so
that only tokens recorded at a prior $\OPwithdraw$ can be spent or redeemed.
Theft prevention is the stronger check that $\algo{tokens}[\vec{\rho}]$ names
the calling user. Double spending detection is the test on
$|\algo{deposited}[\vec{\rho}[0]]|$ in $\OPdeposit$, and identification is
$\OPidentify$ returning the owner recorded at the divergence prefix. That this
owner is recorded rather than inferred is what gives non-frameability, since an
honest user is only ever written into $\algo{tokens}$ as the owner of a token
they were actually paid. Conservation of value is the pairing of the debits in
$\OPwithdraw$ with the credits in $\OPdeposit$, both of the same $v$, with no
interface that alters $\algo{accounts}$ or $\algo{reserves}$ otherwise.
Anonymity and unlinkability are what $\simu$ does not receive: on a withdrawal
by an honest user it learns the user and the value but not $\rho$, and on the
matching deposit it learns $\vec{\rho}$ but has no way to associate it with that
withdrawal.

\begin{figure} \centering \hspace*{-0.05\textwidth}
\begin{minipage}{1.1\textwidth}

    \begin{framed} \scriptsize

      \setlist[enumerate]{ topsep=2pt, partopsep=0pt, parsep=0pt, itemsep=1pt }

      \begin{multicols}{2}

	Before any call to the offline payment interfaces, $\simu$ issues
corruption queries:

	{$\OPcorrupt$.} Upon an active corruption request $(\OPcorrupt, P)$ for
$P \in \algo{Dev} \cup \algo{SE}$ from $\simu$ do: $\algo{mal} \gets \algo{mal}
\cup P$.

	{$\OPcreate$.} Upon receiving for the first time $(\OPcreate,
\algo{usr}, v)$ from $\algo{Int} \in \algo{I}$ for $\algo{usr} \in \algo{usrs}$,
send $(\OPcreate, \algo{usr}, \algo{Int}, v)$ to $\simu$ and wait for the
go-ahead. On receiving it, do:

	\begin{enumerate}

	  \item If $\algo{accounts}[\algo{usr}] \neq \bot$, then stop.

	  \item Otherwise:

	    \begin{enumerate}

	      \item Set $\algo{accounts}[\algo{usr}] \leftarrow \langle
\algo{Int}, v\rangle$.

	      \item Set $\algo{reserves}[\algo{Int}] \leftarrow
\algo{reserves}[\algo{Int}]+v$.

	      \item Return $\OPcreateend$ to $\algo{Int}$.

	    \end{enumerate}

	\end{enumerate}

	{$\OPregister$.} Upon receiving for the first time $(\OPregister,
\algo{usr})$ from $\reg$ for $\algo{usr} \in \algo{usrs}$ that already has an
account, send $(\OPregister, \algo{usr})$ to $\simu$ and wait for the go-ahead.
On receiving it, do:

	\begin{enumerate}

	  \item Set $\algo{registered} \gets \algo{registered} \cup \algo{usr}$.

	  \item Return $\OPregisterend$ to $\reg$.

	\end{enumerate}

	{$\OPwithdraw$.} On receiving a $(\OPwithdraw, \algo{usr} = (\algo{se},
\algo{dev}), \algo{Int}, v)$ from registered user $\algo{usr}$ that has an
account with $\algo{Int}$, send $(\OPwithdraw, \algo{usr}, \algo{Int}, v)$ to
$\algo{Int}$ and $(\OPwithdraw, \algo{Int}, v)$ to $\algo{cb}$ and await their
go-ahead. Then send message $\OPwithdraw$ to $\simu$ and wait for the go-ahead.
On receiving it, fetch $\algo{accounts}[\algo{usr}] \rightarrow \langle
\algo{Int}, v^*\rangle$, and do:

	\begin{enumerate}

	  \item If $v^* < v$ or  $\algo{reserves}[\algo{Int}] < v$, then stop.

	  \item Otherwise, do:

	    \begin{enumerate}

	      \item Set $\algo{accounts}[\algo{usr}] \leftarrow \langle
\algo{Int}, v^* - v\rangle$.

	      \item Set $\algo{reserves}[\algo{Int}] \leftarrow
\algo{reserves}[\algo{Int}] - v$.

	      \item If $\algo{se} \not\in \algo{mal}$, then randomly select
$\rho$.

	      \item Otherwise, receive $\rho$ from $\simu$.

	      \item Write $\vec{\rho} = (\rho)$. If $\algo{tokens}[\vec{\rho}]
\neq \bot$, then stop.

	      \item Otherwise, do:

		\begin{enumerate}

		  \item Set $\algo{tokens}[\vec{\rho}] \gets \langle \algo{usr},
v \rangle$.

		  \item If $\algo{se} \in \algo{mal}$ or $\algo{dev} \in
\algo{mal}$, then return $(\OPwithdraw, \algo{usr}, \rho, v)$ to $\simu$.

		  \item Otherwise, return $(\OPwithdraw, \algo{usr}, v)$ to
$\simu$.

		  \item Return $(\OPwithdraw, \rho, v, \OPwithdrawend)$ to
$\algo{usr}$.

		\end{enumerate}

	    \end{enumerate}

	\end{enumerate}

	{$\OPexecute$.} On receiving $(\OPexecute, \vec{\rho}, v, \algo{usr}' =
(\algo{se}', \algo{dev}'))$ from registered user $\algo{usr} = (\algo{se},
\algo{dev})$ for registered user $\algo{usr}'$, send $(\OPexecute, \vec{\rho},
v)$ to $\algo{usr}'$. On receiving the response $\algo{accept}$ from
$\algo{usr}'$, send $(\OPexecute, \vec{\rho}, v)$ to $\simu$ and wait for the
go-ahead. On receiving it, do:

	\begin{enumerate}

	  \item If $\algo{se}' \not\in \algo{mal}$ and $\algo{dev}' \not\in
\algo{mal}$, then do:

	    \begin{enumerate}

	      \item If $\algo{tokens}[\vec{\rho}] \neq \langle \algo{usr}, v
\rangle$, then stop.

	      \item Otherwise, do:

		\begin{enumerate}

		  \item Randomly select a unique random number $\rho'$.

		  \item Set $\algo{tokens}[\vec{\rho'}] \gets
\langle\algo{usr}', v \rangle$ where $\vec{\rho'} = (\vec{\rho}, \rho')$.

		  \item If $\algo{se} \in \algo{mal}$ or $\algo{dev} \in
\algo{mal}$, then return $(\OPexecute, \vec{\rho'}, v, \algo{usr}, \algo{usr}')$
to $\simu$.

		  \item Otherwise, return $(\OPexecute, \vec{\rho'}, v)$ to
$\simu$.

		\end{enumerate}

	    \end{enumerate}

	  \item Otherwise, do:

	    \begin{enumerate}

	      \item If $\algo{se}' \in \algo{mal}$, then receive $\rho'$ from
$\simu$.

	      \item Otherwise, randomly select $\rho'$.

	      \item Set $\vec{\rho'} = (\vec{\rho}, \rho')$.

	      \item If $\algo{tokens}[\vec{\rho}] = \langle \algo{usr}, v
\rangle$, then set $\algo{tokens}[\vec{\rho'}] \leftarrow \langle \algo{usr}', v
\rangle$. A corrupt payee that accepts a token the payer does not own records
nothing, so cannot pass it on.

	      \item Send $(\OPexecute,\vec{\rho'}, v, \algo{usr}, \algo{usr}')$
to $\simu$.

	    \end{enumerate}

	  \item Return $(\OPexecute, \vec{\rho'}, v, \OPexecuteend)$ to
$\algo{usr}'$, where $\vec{\rho'}$ is as set above. If branch 1 stopped, return
nothing.

	\end{enumerate}

	{$\OPdeposit$.} On receiving $(\OPdeposit, \vec{\rho}, v, \algo{Int})$
from registered user $\algo{usr}$ that has an account with $\algo{Int}$, send
$(\OPdeposit, \vec{\rho}, v, \algo{Int}, \algo{usr})$ to $\algo{Int}$ and
$(\OPdeposit, v, \algo{Int})$ to $\algo{cb}$ and await their go-ahead. Then send
$(\OPdeposit, \vec{\rho}, v, \algo{Int}, \algo{usr})$ to $\simu$ and wait for
the go-ahead. On receiving it, do:

	\begin{enumerate}

	  \item If $\algo{tokens}[\vec{\rho}] \neq \langle \algo{usr}, v
\rangle$, then stop.

	  \item Otherwise, do:

	    \begin{enumerate}

	      \item Set $\algo{deposited}[\vec{\rho}[0]] \gets
\algo{deposited}[\vec{\rho}[0]] \cup \{\vec{\rho}\}$.

	      \item If $|\algo{deposited}[\vec{\rho}[0]]| > 1$, then stop.

	      \item Otherwise, do:

		\begin{enumerate}

		  \item Fetch $\algo{accounts}[\algo{usr}] \rightarrow \langle
\algo{Int}, v^*\rangle$.

		  \item Set $\algo{accounts}[\algo{usr}] \gets \langle
\algo{Int}, v^*+v \rangle$.

		  \item Set $\algo{reserves}[\algo{Int}] \gets
\algo{reserves}[\algo{Int}] + v$.

		  \item Return $(\OPdeposit, \OPdepositend)$ to $\algo{usr}$.

		\end{enumerate}

	    \end{enumerate}

	\end{enumerate}

	{$\OPidentify$.} On receiving $(\OPidentify, \rho)$ from $\auditor$,
send $(\OPidentify, \rho)$ to $\simu$ and wait for the go-ahead. On receiving
it, do:

	\begin{enumerate}

	  \item If $|\algo{deposited}[\rho]| \leq 1$, then return $(\OPidentify,
\emptyset, \OPidentifyend)$ to $\auditor$.

	  \item Otherwise, do:

	    \begin{enumerate}

	      \item Fetch $\algo{deposited}[\rho] \rightarrow \{\vec{\rho}_1,
\cdots, \vec{\rho}_n\}$.

	      \item Set $\algo{DS} \leftarrow \emptyset$.

	      \item $\forall i\neq j$, do:

		\begin{enumerate}

		  \item Let $\vec{\rho}_{ij}$ denote the longest common prefix
of $\vec{\rho}_i$ and $\vec{\rho}_j$, that is $\vec{\rho}_{ij} = (\rho_0,
\ldots, \rho_{k-1})$ with $\vec{\rho}_i[l] = \vec{\rho}_j[l]$ for all $l < k$
and either $\vec{\rho}_i[k] \neq \vec{\rho}_j[k]$ or one of the two has length
$k$.

		  \item Fetch $\algo{tokens}[\vec{\rho}_{ij}] \rightarrow
\langle \algo{usr}_{ij}, v\rangle$.

		  \item Set $\algo{DS} \gets \algo{DS} \cup \algo{usr}_{ij}$.

		\end{enumerate}

	      \item Return $(\OPidentify, \algo{DS}, \OPidentifyend)$ to
$\auditor$ and $\algo{cb}$.

	    \end{enumerate}

	\end{enumerate}

      \end{multicols}

    \end{framed}

  \end{minipage}

\caption{Ideal functionality $\idfu{\algo{op}}$ for offline payments.}
\label{fig:ideal-func} \end{figure}
